\section{Электростатическое кольцо}
\label{sec:electrostatic}

В главах~\ref{sec:magnetic} и~\ref{sec:ch2} модуль скорости между зазорами зависит только от полного импульса. Это свойство статического магнитного поля, и из него следует вынос $v$ из интеграла~\eqref{eq:mag-T}. В электростатическом кольце кинетическая энергия вдоль орбиты меняется вместе с потенциалом, и $v$ остаётся под интегралом. Уравнение энергии на зазоре при этом то же, что~\eqref{eq:mag-energy}: высокочастотное поле видит опоздание $T-T_s$ и за один пролёт меняет полную энергию на $eV\sin\varphi$. Отличие электростатики сидит в периоде $T=\oint\mathrm{d}l/v$, а не в приращении энергии.

\subsection{Переменная импульса}
\label{sec:es-delta}

Между зазорами сохраняется полная энергия
\begin{equation}
\label{eq:es-energy}
E=\gamma mc^2+q\Phi.
\end{equation}
Потенциал на проектной дуге выбран так, что на опорном радиусе $\Phi=0$, и в прямой $\Phi=0$. Величина $\delta$ этой главы есть относительный импульс частицы, у которой та же полная энергия и нулевой потенциал: $p=p_s(1+\delta)$ в зазоре и в прямой. На зазоре $\mathrm{d}E=\beta_s^2 E\,\mathrm{d}\delta$, поэтому уравнение~\eqref{eq:mag-energy} переносится сюда без изменения коэффициентов. Локальный импульс внутри дефлектора этим $\delta$ не является: часть полной энергии сидит в $q\Phi$.

В COSY шестая координата устроена иначе. Процедура \texttt{RP} задаёт кинетическую энергию опорной частицы, и
\begin{equation}
\label{eq:es-deltak}
\delta_K=\frac{K-K_0}{K_0}
\end{equation}
есть относительная кинетическая энергия в опорной точке, где потенциал равен опорному. В статическом поле эта координата постоянна. На нулевом потенциале она связана с $\delta$ рядом
\begin{equation}
\label{eq:es-dk-series}
\delta_K
=\frac{\gamma_s+1}{\gamma_s}\,\delta
+\frac{\gamma_s+1}{2\gamma_s^3}\,\delta^2
+O(\delta^3).
\end{equation}
Пятая координата $l$ есть не приращение пути. Её производная по дуге в уравнениях COSY пропорциональна отличию локального пролётного множителя от единицы, и на целом обороте
\begin{equation}
\label{eq:es-cosy-T}
\frac{T-T_s}{T_s}
=-\frac{\gamma_s+1}{\gamma_s}\,\frac{l}{C_s},
\end{equation}
пока на опорной орбите $C_s=v_s T_s$. На чистом дрейфе, где $\alpha_0=\alpha_1=0$, подстановка~\eqref{eq:es-dk-series} в~\eqref{eq:es-cosy-T} возвращает магнитный ряд~\eqref{eq:mag-vsv}: коэффициент при $\delta$ равен $-1/\gamma_s^2$, коэффициент при $\delta^2$ равен $1/\gamma_s^4+3\beta_s^2/(2\gamma_s^2)$. Электрический потенциал в эту проверку не входил. В кольце с дефлекторами он уже сидит внутри отображения $l$, потому что локальная кинетическая энергия в уравнениях COSY уменьшается на $q\Phi$. Второй раз применять к $l$ магнитные формулы~\eqref{eq:mag-eta0} и~\eqref{eq:mag-eta1} не нужно.

Коэффициенты Тейлора COSY понимаются как множители мономов: $l\supset(5|6)\,\delta_K+(5|66)\,\delta_K^2$. Вместе с~\eqref{eq:es-dk-series} и~\eqref{eq:es-cosy-T} они дают
\begin{align}
\eta_0
&=
-\left(\frac{\gamma_s+1}{\gamma_s}\right)^2\frac{(5|6)}{C_s},
\label{eq:es-eta0-map}\\
\eta_1
&=
-\frac{\gamma_s+1}{\gamma_s C_s}
\Biggl(
(5|6)\,\frac{\gamma_s+1}{2\gamma_s^3}
+(5|66)\left(\frac{\gamma_s+1}{\gamma_s}\right)^2
\Biggr).
\label{eq:es-eta1-map}
\end{align}
Это $\eta_0$ и $\eta_1$ периода, электростатические или магнитные, смотря по тому, какое поле породило отображение.

\subsection{Линейный слип-фактор}
\label{sec:es-eta0}

Период по-прежнему есть~\eqref{eq:mag-T}. Опорная скорость $v_s$ отвечает $\gamma_s$ и $\Phi=0$. Относительное опоздание
\begin{equation}
\label{eq:es-rel}
\frac{T-T_s}{T_s}
=\frac{1}{C_s}\int\Biggl(\frac{v_s}{v}\,\frac{\mathrm{d}l}{\mathrm{d}s}-1\Biggr)\mathrm{d}s.
\end{equation}
В прямой $\Phi=0$, локальный лоренц-фактор совпадает с $\gamma(\delta)$ формулы~\eqref{eq:mag-gamma-delta}, и $v_s/v$ есть ряд~\eqref{eq:mag-vsv}. В цилиндрическом дефлекторе с полем, спадающим как $1/r$,
\begin{equation}
\label{eq:es-phi}
\frac{q\Phi}{mc^2}
=\beta_s^2\gamma_s\ln\Bigl(1+\frac{x}{\rho}\Bigr)
=\beta_s^2\gamma_s\Biggl(\frac{x}{\rho}-\frac{1}{2}\Biggl(\frac{x}{\rho}\Biggr)^2+\cdots\Biggr).
\end{equation}
Коэффициент $\beta_s^2\gamma_s/\rho$ есть электрическое поле, которое держит опорную частицу на радиусе $\rho$. Локальная кинетическая энергия равна $E-q\Phi$.

На дисперсионной орбите $x=D\delta+x_2\delta^2+\cdots$. До первого порядка по $\delta$ избыток кинетического лоренц-фактора в дуге есть $\beta_s^2\gamma_s(\delta-x/\rho)$. Обращение скорости, то же, что приводит к первому члену~\eqref{eq:mag-vsv}, даёт
\begin{equation}
\label{eq:es-v-linear}
\frac{v_s}{v}-1
=-\frac{\delta}{\gamma_s^2}+\frac{1}{\gamma_s^2}\frac{x}{\rho}.
\end{equation}
Элемент пути до того же порядка есть $\mathrm{d}l/\mathrm{d}s=1+x/\rho$. Произведение в~\eqref{eq:es-rel} с точностью до $\delta$ равно
\begin{equation}
\label{eq:es-integrand}
\frac{v_s}{v}\,\frac{\mathrm{d}l}{\mathrm{d}s}-1
=-\frac{\delta}{\gamma_s^2}
+\Biggl(1+\frac{1}{\gamma_s^2}\Biggr)\frac{x}{\rho}
=-\frac{\delta}{\gamma_s^2}
+(2-\beta_s^2)\frac{x}{\rho}.
\end{equation}
Второе равенство есть $1+1/\gamma_s^2=2-\beta_s^2$. В прямой $x/\rho$ не возникает, и остаётся только $-\delta/\gamma_s^2$. Интеграл геометрического члена по кольцу есть $\alpha_0\delta$, если $x=D\delta$ и весь изгиб электрический, с одним и тем же смыслом $\rho$ в пути и в~\eqref{eq:es-phi}. Поэтому
\begin{equation}
\label{eq:es-eta0}
\eta_0^{\mathrm{es}}
=(2-\beta_s^2)\alpha_0-\frac{1}{\gamma_s^2}
=\alpha_0-\frac{1-\alpha_0}{\gamma_s^2}.
\end{equation}
Множитель $(2-\beta_s^2)$ стоит перед всем $\alpha_0$, когда каждая дуга электрическая, опорный потенциал на проектном радиусе равен потенциалу прямой и линейная часть потенциала есть именно~\eqref{eq:es-phi}. Добавка к геометрическому $\alpha_0$ равна $\alpha_0/\gamma_s^2$: на смещённой наружу орбите скорость падает, и $1/v$ растёт пропорционально $x/\rho$.

Если часть изгиба магнитная, геометрический $\alpha_0=(1/C_s)\int D/\rho\,\mathrm{d}s$ собирает полную кривизну, а падение скорости даёт только электрическая доля. Пусть $qE=\beta_s^2\gamma_s mc^2/\rho_E$ определяет электрический радиус, и
\begin{equation}
\label{eq:es-alphaE}
\alpha_0^{E}
=\frac{1}{C_s}\int\frac{D}{\rho_E}\,\mathrm{d}s.
\end{equation}
Тогда
\begin{equation}
\label{eq:es-eta0-gen}
\eta_0
=\alpha_0+\frac{\alpha_0^{E}}{\gamma_s^2}-\frac{1}{\gamma_s^2}.
\end{equation}
У чисто магнитного кольца $\alpha_0^{E}=0$, и~\eqref{eq:es-eta0-gen} совпадает с~\eqref{eq:mag-eta0}. У чисто электрического $\alpha_0^{E}=\alpha_0$, и получается~\eqref{eq:es-eta0}. У фильтра Вина суммарная кривизна опорной орбиты равна нулю, но $\rho_E$ конечен, $\alpha_0^{E}\neq 0$, и линейный слип-фактор не сводится к $-1/\gamma_s^2$.

Формула~\eqref{eq:es-eta0} обращается в нуль при
\begin{equation}
\label{eq:es-transition}
(2-\beta_s^2)\alpha_0\gamma_s^2=1,
\end{equation}
то есть при $\gamma_s^2=1/\alpha_0-1$. Магнитный переход того же $\alpha_0$ остаётся при $\gamma^2=1/\alpha_0$. Электростатический переход лежит ниже: $\sqrt{1/\alpha_0-1}<1/\sqrt{\alpha_0}$. Кольцо с $\gamma_s$ между этими двумя значениями имеет $\eta_0^{\mathrm{es}}$ и магнитное $\eta_0=\alpha_0-1/\gamma_s^2$ разных знаков.

Гладкий цилиндр без прямых даёт крайний случай. На каждой круговой орбите в поле $1/r$ кинетическая энергия равна опорной, скорость равна $v_s$, и период пропорционален радиусу. При $\alpha_0=1$ формула~\eqref{eq:es-eta0} даёт $\eta_0^{\mathrm{es}}=1$, а магнитная формула дала бы $\beta_s^2$. Разложение радиуса по $\delta$ до второго порядка даёт ещё
\begin{equation}
\label{eq:es-smooth-eta1}
\eta_1=\frac{1}{2}+\frac{1}{2\gamma_s^2}.
\end{equation}
То же значение получается и для длины, и для времени: скорость от $\delta$ не зависит.

\subsection{Формула (25) и множитель $F_v$}
\label{sec:es-fv}

В~\cite{Senichev2013} переход к электростатике записан одним множителем $F_v$ при кинематических членах. Линейная строка формулы (25) той работы есть
\begin{equation}
\label{eq:es-sen25}
\frac{\Delta T}{T}
=\Biggl(\alpha_0-\frac{1}{F_v\gamma_s^2}\Biggr)\delta.
\end{equation}
Сравнение с~\eqref{eq:es-eta0} фиксирует
\begin{equation}
\label{eq:es-fv}
F_v=\frac{1}{1-\alpha_0}.
\end{equation}
Это значение не зависит от энергии. Оно воспроизводит линейный слип-фактор целиком, включая сдвиг перехода~\eqref{eq:es-transition}. Условие смены знака $F_v>1/(\alpha_0\gamma_s^2)$, выписанное в той работе, при~\eqref{eq:es-fv} совпадает с $\gamma_s^2>1/\alpha_0-1$. Отдельного механизма смены знака, помимо сдвига критической энергии, линейный член не содержит. Для гладкого цилиндра $\alpha_0=1$ и $F_v$ расходится: скорость на семействе круговых орбит от $\delta$ не зависит, что согласовано с $\eta_0^{\mathrm{es}}=1$.

Тот же $F_v$ не переносит квадратичную скобку формулы (25). В ней по-прежнему нет $3\beta_s^2/(2\gamma_s^2)$, уже отсутствовавшего в магнитной записи, и нет членов от кривизны потенциала $\varphi_2$ и от $\langle(D/\rho)^2\rangle$. Подстановка~\eqref{eq:es-fv} в скобку $\alpha_1-\alpha_0/(F_v\gamma_s^2)+1/(F_v^2\gamma_s^4)$ на модельном кольце раздела~\ref{sec:es-num} даёт $0{,}432$ против измеренного $\eta_1^{\mathrm{es}}=0{,}759$. Если к этой скобке вручную прибавить $3\beta_s^2/(2\gamma_s^2)$, получается $0{,}776$: ближе, но не то же число. Бетатронное слагаемое $(\Delta L/L)_\beta$ в (25) оставлено с коэффициентом $1$. Это коэффициент пути. В период входит время.

\subsection{Второй порядок}
\label{sec:es-eta1}

Квадратичный коэффициент периода получается тем же произведением $(v_s/v)(\mathrm{d}l/\mathrm{d}s)$, если удержать $\delta^2$ в~\eqref{eq:es-phi}, в $\gamma(\delta)$ и в $x=D\delta+x_2\delta^2$, а в элементе пути ещё $x'^2/2$. Пусть $d_1=D/\rho$, $d_2=x_2/\rho$ и угловые скобки есть $(1/C_s)\int(\cdots)\,\mathrm{d}s$. Геометрические $\alpha_0=\langle d_1\rangle$ и $\alpha_1$ сохраняют смысл главы~\ref{sec:magnetic}: $\alpha_1$ включает и $\langle d_2\rangle$, и $\langle D'^2\rangle/2$. Магнитный коэффициент~\eqref{eq:mag-eta1} на тех же $\alpha_0$ и $\alpha_1$ обозначается $\eta_1^{\mathrm{mag}}$. Прямое разложение даёт
\begin{equation}
\label{eq:es-eta1}
\eta_1^{\mathrm{es}}
=\eta_1^{\mathrm{mag}}
+\frac{1}{\gamma_s^2}
\Biggl(
\langle d_2\rangle
+2\Biggl\langle\Biggl(\frac{D}{\rho}\Biggr)^2\Biggr\rangle
-3\alpha_0
\Biggr).
\end{equation}
Добавка собрана из потенциала. В неё входят квадрат линейной дисперсии и вторая орбита $x_2$; одного $\alpha_1$ недостаточно, потому что $\langle D'^2\rangle$ и $\langle d_2\rangle$ входят в $\alpha_1$ одной суммой, а в добавку к периоду идёт только $\langle d_2\rangle$. На гладком цилиндре $d_1=1$, $D'=0$ и $\langle d_2\rangle=\eta_1$ формулы~\eqref{eq:es-smooth-eta1}. Подстановка обращает~\eqref{eq:es-eta1} в тождество с этим $\eta_1$.

Секступоль, как и в магните, меняет $\alpha_1$ и вместе с ним $\eta_1^{\mathrm{es}}$. Обнуление магнитной комбинации $\eta_1^{\mathrm{mag}}$ оставляет добавку~\eqref{eq:es-eta1}.

\subsection{Бетатронное опоздание}
\label{sec:es-lambda}

При $\delta=0$ полная энергия равна $E_s$, а локальная скорость в дуге зависит от $x$. Относительное опоздание одного оборота после среднего по бетатронной фазе записывается суммой
\begin{equation}
\label{eq:es-LT}
\Lambda_T=\Lambda_C+\Lambda_\Phi.
\end{equation}
$\Lambda_C$ есть геометрическое удлинение главы~\ref{sec:ch2}: среднее от $x_{\mathrm{nl}}/\rho$ и от $(x_\beta'^2+y_\beta'^2)/2$, делённое на $C_s$. В магните $\Lambda_\Phi=0$, потому что $|v|$ от $x$ не зависит, и $\Lambda_T=\Lambda_C$. В электрической дуге при $\delta=0$ ряд~\eqref{eq:es-phi} даёт
\begin{equation}
\label{eq:es-v-beta}
\frac{v_s}{v}-1
=\frac{1}{\gamma_s^2}\frac{x}{\rho}
+\frac{1}{\gamma_s^2}\Biggl(\frac{x}{\rho}\Biggr)^2
+O(x^3).
\end{equation}
Линейный член после среднего оставляет $\langle x_{\mathrm{nl}}\rangle/(\rho\gamma_s^2)$. Квадрат $x_\beta$ оставляет $\langle x_\beta^2\rangle/(\rho^2\gamma_s^2)$. Произведение линейного падения скорости на линейное удлинение пути даёт ещё один такой же квадрат. Угловой член $x'^2/2$ потенциала не меняет и входит только в $\Lambda_C$. Поэтому на цилиндрических дугах
\begin{equation}
\label{eq:es-Lphi}
\Lambda_\Phi
=\frac{1}{C_s}
\int_{E}
\Biggl(
\frac{\langle x_{\mathrm{nl}}\rangle}{\rho\gamma_s^2}
+\frac{2\langle x_\beta^2\rangle}{\rho^2\gamma_s^2}
\Biggr)\mathrm{d}s,
\end{equation}
где $\langle x_\beta^2\rangle=J_x\beta_x$ и интеграл идёт по электрическим дугам. Сдвиг средней орбиты входит в время с тем же весом, что $x_{\mathrm{nl}}/\rho$ входит в путь: в $\Lambda_T$ перед $\langle x_{\mathrm{nl}}\rangle/\rho$ стоит $2-\beta_s^2$, перед квадратом угла стоит единица, а член $\langle x_\beta^2\rangle/\rho^2$ есть только у потенциала. В прямой $\Lambda_\Phi$ не возникает.

\subsection{Связь с хроматичностью}
\label{sec:es-xi}

Каноническая структура однооборотного отображения связывает среднее время с производной тюна. Пусть $s$ --- независимая переменная, $J_{\mathrm{can}}=(1/2\pi)\oint p_x\,\mathrm{d}x$ --- действие при каноническом $p_x$ в единицах $mc$, а полная энергия $E$ в единицах $mc^2$ сопряжена времени так, что импульс времени есть $-E$. Равенство смешанных производных гамильтониана, порождающего сдвиг по $s$, даёт
\begin{equation}
\label{eq:es-dTdJ}
\frac{\partial\langle T\rangle}{\partial J_{\mathrm{can}}}
=-2\pi\frac{\partial\nu}{\partial E}.
\end{equation}
На прямой при $\Phi=0$ и $\delta=0$ канонический импульс и геометрический угол связаны множителем $p_s=\beta_s\gamma_s$, так что $J_{\mathrm{can}}=p_s J$ для твисс-действия $J$ пары $(x,x')$. Энергия в зазоре и переменная $\delta$ связаны $\mathrm{d}E/\mathrm{d}\delta=\beta_s^2 E$. При $\xi=\partial\nu/\partial\delta$ отсюда
\begin{equation}
\label{eq:es-sym}
\frac{1}{T_s}\frac{\partial\langle T\rangle}{\partial J}
=-\frac{2\pi\xi}{v_s T_s}.
\end{equation}
На опорной орбите обеих моделей этого раздела кинетическая энергия всюду равна опорной, $v_s T_s=C_s$, и правая часть равна $-2\pi\xi/C_s$. В магните то же число есть производная пути: $|v|$ от $x$ не зависит, и $\Lambda_C$ формулы~\eqref{eq:beta-xi} совпадает с $\Lambda_T$. В электростатике~\eqref{eq:es-sym} относится ко времени. Производная пути к $\xi$ этой формулой не привязана.

Модельное кольцо раздела~\ref{sec:es-num} этой нормы не воспроизводит. Фазовое среднее одного оборота на линейном эллипсе Твисса, при действии от $10^{-8}$ до $4\times 10^{-5}\,\mathrm{м}$, даёт для обеих машин прямую
\begin{equation}
\label{eq:es-scan}
\frac{1}{J}\Lambda
=\kappa\Biggl(-\frac{2\pi\xi}{C_s}\Biggr)+b,
\end{equation}
где $\Lambda$ есть $\Lambda_T$ или $\Lambda_C$. Коэффициенты одни на весь скан тонкого секступоля.

В магните время и путь совпадают. Для них $\kappa=1{,}076$ и $b=1{,}00\times 10^{-3}\,\mathrm{м}^{-1}$. При естественной хроматичности $\xi=-0{,}207$ предсказание $-2\pi\xi/C_s$ равно $6{,}85\times 10^{-3}\,\mathrm{м}^{-1}$, измеренный наклон равен $8{,}36\times 10^{-3}\,\mathrm{м}^{-1}$, и отношение равно $1{,}22$. Та же прямая при $S=0{,}002$ даёт отношение $0{,}95$, при $S=0{,}004$ --- отношение $1{,}03$. Отношение $1{,}22$ есть значение~\eqref{eq:es-scan} при естественном $\xi$, а не отдельный дефект точки $S=0$. Аддитивный член $b$ при $S=0{,}001$, где $\xi\approx 0{,}012$ и $-2\pi\xi/C_s$ уже меньше $b$, переворачивает знак наклона относительно $-2\pi\xi/C_s$: измеренный наклон времени остаётся около $+5\times 10^{-4}\,\mathrm{м}^{-1}$.

Угловая часть пути на этой оптике равна $(1/(2C_s))\int\gamma_x\,\mathrm{d}s=2{,}41\times 10^{-2}\,\mathrm{м}^{-1}$. Измеренный путь $8{,}36\times 10^{-3}\,\mathrm{м}^{-1}$ меньше неё втрое: средний сдвиг орбиты снимает большую долю квадрата угла. Остаток после этого сокращения и есть величина, которую~\eqref{eq:es-scan} сравнивает с $-2\pi\xi/C_s$. Почему коэффициент при секступольном скане равен $1{,}076$, а не единице, и откуда аддитивные $1{,}00\times 10^{-3}\,\mathrm{м}^{-1}$, этот расчёт не закрывает. Определитель линейной матрицы $(x,x')$ равен единице с точностью $10^{-6}$, действие согласовано с Твиссом, и оба наклона не меняются при увеличении $J$ в четыре тысячи раз. Расхождение сидит в соответствии жёсткого однородного сектора и тонкой линзы тождеству Шоджи после взаимного сокращения $\int\gamma\,\mathrm{d}s$ и $\langle x_{\mathrm{nl}}\rangle$. Это открытый вопрос контроля модели. На фоне электростатического сдвига ниже он мал: $1\times 10^{-3}$ против $1\times 10^{-2}$.

В электростатическом кольце время и путь расходятся. По времени $\kappa_T=1{,}227$ и $b_T=-9{,}65\times 10^{-3}\,\mathrm{м}^{-1}$. По пути $\kappa_C=0{,}747$ и $b_C=-2{,}05\times 10^{-2}\,\mathrm{м}^{-1}$. Среднеквадратичный остаток прямой на скане $S=0\ldots 0{,}012$ меньше $3\times 10^{-5}\,\mathrm{м}^{-1}$. При $\xi=0$ наклон времени равен $b_T$, а не нулю. Наклон пути равен $b_C$ и тоже не нуль. Разность $b_T-b_C=1{,}09\times 10^{-2}\,\mathrm{м}^{-1}$ есть $\Lambda_\Phi/J$ в той точке, где хроматичность уже погашена одним семейством секступолей. Гашение $\xi$ не гасит опоздание.

Формула~\eqref{eq:es-sym} остаётся следствием гамильтоновой структуры для той пары переменных, в которой она выведена. Жёсткая модель к коэффициенту $1$ и к нулевому $b$ не приходит. Для равновесия ниже используется измеряемое $\Lambda_T$, а не подстановка $-2\pi\xi J/C_s$.

Отсюда же читается ответ про три семейства. В магните обращение $\Lambda$ в нуль вместе с $\xi_x$ и $\xi_y$ есть содержание~\eqref{eq:beta-xi}, с оговоркой предыдущего абзаца о коэффициенте $1{,}076$ и малом $b$. В электростатике условие $\langle\Delta T/T_s\rangle=0$ требует $\eta_1^{\mathrm{es}}=0$ и $\Lambda_T=0$ в каждой плоскости. Это по-прежнему три числа, и трёх семейств хватает, если матрица их отклика невырождена. Эти три числа не совпадают с $\eta_1^{\mathrm{mag}}$, $\xi_x$ и $\xi_y$. Хроматичность, обращённая в нуль, оставляет $b_T$. Читать $\Lambda_T$ нужно из временной координаты отображения. На эллипсе Твисса
\begin{equation}
\label{eq:es-map-L}
\langle l\rangle
=\bigl[(5|11)\beta_x-(5|12)\alpha_x+(5|22)\gamma_x\bigr]J_x
+\bigl[(5|33)\beta_y-(5|34)\alpha_y+(5|44)\gamma_y\bigr]J_y,
\end{equation}
и $\Lambda_T$ получается из~\eqref{eq:es-cosy-T}. Производная тюна по $\delta_K$, которую даёт процедура хроматичности, переводится в $\xi=\partial\nu/\partial\delta$ множителем $(\gamma_s+1)/\gamma_s$ и в~\eqref{eq:es-map-L} сама по себе не подставляется.

На гладком цилиндре линейный тюн равен $\sqrt{2-\beta_s^2}$ на круговой орбите любого радиуса: кинетическая энергия на этих орбитах одна и та же, а набег фазы за оборот от радиуса не зависит. Линейная хроматичность равна нулю. Численно отделить коэффициент при $J$ от линейного по амплитуде колебания периода на коротком слежении не удалось; в приложение это значение не входит.

\subsection{Равновесие}
\label{sec:es-eq}

Подстановка периода в уравнение фазы главы~\ref{sec:magnetic} повторяет вывод раздела~\ref{sec:ch2-eq}. На синхротронном времени $\Lambda_T$ постоянно, $\langle\delta^2\rangle=\delta_m^2/2$, и среднее опоздание обращается в нуль при
\begin{equation}
\label{eq:es-deq}
\delta_{\mathrm{eq}}
=-\frac{1}{\eta_0^{\mathrm{es}}}
\Biggl(
\frac{\eta_1^{\mathrm{es}}\,\delta_m^2}{2}
+\Lambda_T
\Biggr).
\end{equation}
Здесь $\eta_0^{\mathrm{es}}$ и $\eta_1^{\mathrm{es}}$ берутся из времени, формулы~\eqref{eq:es-eta0} и~\eqref{eq:es-eta1} или прямо из~\eqref{eq:es-eta0-map} и~\eqref{eq:es-eta1-map}, а $\Lambda_T$ --- из~\eqref{eq:es-map-L}. Условие, при котором сдвиг исчезает при любой амплитуде синхротронных колебаний и при любых действиях, есть
\begin{equation}
\label{eq:es-comp}
\eta_1^{\mathrm{es}}=0,
\qquad
\Lambda_T=0.
\end{equation}
Второе равенство не заменяется на $\xi_x=\xi_y=0$. Частота при $\cos\varphi=1$ сохраняет вид~\eqref{eq:beta-freq}, если в неё подставить эти $\eta_1^{\mathrm{es}}$ и $\Lambda_T$.

В комбинированном кольце линейный знаменатель есть~\eqref{eq:es-eta0-gen}. Квадратичный коэффициент и $\Lambda_T$ по-прежнему читаются из временной координаты; доля электрического изгиба сама по себе их не собирает.

\subsection{Численная проверка}
\label{sec:es-num}

Модель --- плоское кольцо из восьми ячеек. Протон, кинетическая энергия $233\,\mathrm{МэВ}$, $\gamma_s=1{,}248329$, $\beta_s^2=0{,}358285$. Дуга --- цилиндрический дефлектор радиуса $20\,\mathrm{м}$ на четверть оборота либо однородный магнит того же радиуса и угла; прямые по $8\,\mathrm{м}$; тонкие квадруполи с силой $\pm 0{,}02\,\mathrm{м}^{-1}$ и, в скане, тонкий секступоль на фокусирующих плоскостях. Длина $C_s=189{,}664\,\mathrm{м}$. Потенциал цилиндра равен~\eqref{eq:es-phi} внутри дуги и нулю снаружи; на входе и выходе продольный импульс получает скачок, который сохраняет полную энергию. Магнитное кольцо интегрируется в том же коде с тем же шагом. Подробности и таблица скана --- в приложении~\ref{app:ch3track}.

Линейный коэффициент времени совпадает с формулой с остатком подгонки $10^{-13}$. При $\alpha_0=0{,}637450$
\begin{equation}
\label{eq:es-num-eta0}
\eta_0^{\mathrm{es}}=0{,}404796,
\qquad
\eta_0^{\mathrm{mag}}=\alpha_0-\frac{1}{\gamma_s^2}=-0{,}004265.
\end{equation}
Кольцо лежит ниже магнитного перехода ($\gamma_{\mathrm{tr}}=1/\sqrt{\alpha_0}=1{,}253$) и выше электростатического ($\gamma_{\mathrm{tr}}=\sqrt{1/\alpha_0-1}=0{,}754$). Знаки $\eta_0$ поэтому различны. Это и есть смена знака слип-фактора при неизменных $\alpha_0$ и $\gamma_s$: критическая энергия сдвинулась, а рабочая точка осталась на месте.

Квадратичный коэффициент при нулевом секступоле равен $0{,}75898$. Магнитная формула на измеренных $\alpha_0$ и $\alpha_1=0{,}52578$ даёт $0{,}87339$. Формула~\eqref{eq:es-eta1}, в которой $\langle(D/\rho)^2\rangle$ и $\langle d_2\rangle$ взяты с замкнутых орбит, после вычитания члена $\delta^4$ в $x_2$ отличается от $0{,}75898$ на $1\times 10^{-4}$.

Для решётки \texttt{electrostatic\_mh} паспортные значения --- протон $233\,\mathrm{МэВ}$, $C=419{,}1\,\mathrm{м}$, пятьдесят дуг по $7{,}2^\circ$ и $4{,}685\,\mathrm{м}$, $\alpha_0=0{,}0964$, естественные хроматичности $\xi_x=-9{,}17$ и $\xi_y=-8{,}82$. При поле $112{,}5\,\mathrm{кВ/см}$ электрическая жёсткость $\beta_s^2\gamma_s mc^2/(qE)$ даёт радиус $37{,}3\,\mathrm{м}$, и длина дуги, делённая на этот радиус, совпадает с $7{,}2^\circ$ в пределах округления поля. На этих $\alpha_0$ и $\gamma_s$
\begin{equation}
\label{eq:es-mh}
\eta_0^{\mathrm{es}}=-0{,}48345,
\qquad
\eta_0^{\mathrm{mag}}=-0{,}54531.
\end{equation}
Оба значения отрицательны: $\gamma_{\mathrm{tr}}^{\mathrm{es}}=3{,}062$ и $\gamma_{\mathrm{tr}}^{\mathrm{mag}}=3{,}221$ лежат выше $\gamma_s=1{,}248$. Паспортное $\eta_1=1{,}07$ получено магнитной формулой~\eqref{eq:mag-eta1} и коэффициентом $\eta_1^{\mathrm{es}}$ не является. Его даёт~\eqref{eq:es-eta1-map}. Бетатронное $\Lambda_T$ этой решётки из хроматичностей не собирается: в модели раздел~\ref{sec:es-xi} при $\xi=0$ остаётся наклоном $-9{,}65\times 10^{-3}\,\mathrm{м}^{-1}$.
